% =====================================================================
%  Redes Neuronales 2026 (FaMAF-UNC) - Trabajo Práctico 1
%  Boceto de informe: "El modelo de neuronas de Izhikevich"
%  Actualizado a la consigna del 28/9/2026 (notación a, b, c, d del artículo,
%  corrientes y condiciones iniciales de la consigna, nueva segunda parte).
%
%  Cómo usarlo:
%   - Overleaf: New Project > Upload Project > subí el .zip con main.tex y
%     la carpeta figuras/. Compilá con pdfLaTeX (el botón "Recompile").
%   - Local: latexmk -pdf main.tex   (o pdflatex main.tex dos veces).
%
%  Todo lo que está entre [corchetes] o marcado con TODO es para revisar o
%  reemplazar. Los números del texto salen de tp1-izhikevich/codigo/resultados.json
%  (semilla de la red: 2026): verifíquenlos contra sus propias corridas.
%  Límite de la cátedra: el PDF final NO debe superar las 4 páginas.
% =====================================================================
\documentclass[10pt,a4paper,twocolumn]{article}

% ---------- Idioma y codificación ----------
\usepackage[utf8]{inputenc}
\usepackage[T1]{fontenc}
\usepackage[spanish,es-tabla,es-noshorthands]{babel} % "Tabla" en vez de "Cuadro"
\usepackage{lmodern}

% ---------- Página ----------
\usepackage[a4paper,margin=1.8cm,columnsep=0.7cm]{geometry}
\usepackage{microtype}

% ---------- Matemática y unidades ----------
\usepackage{amsmath,amssymb}
\usepackage{siunitx}
\sisetup{output-decimal-marker={,}}          % 0,04 en vez de 0.04
\DeclareSIUnit{\milivolt}{mV}

% ---------- Figuras y tablas ----------
\usepackage{graphicx}
\graphicspath{{figuras/}}                     % las imágenes van en figuras/
\usepackage{subcaption}                       % paneles (a), (b) dentro de una figura
\usepackage{booktabs}                         % líneas lindas en tablas
\usepackage[font=small,labelfont=bf]{caption}

% ---------- Enlaces ----------
\usepackage{xcolor}
\usepackage[colorlinks=true,linkcolor=blue!50!black,citecolor=blue!50!black,urlcolor=blue!50!black]{hyperref}

% ---------------------------------------------------------------------
% \figura{archivo}{ancho}: incluye la imagen si existe en figuras/; si
% todavía no la generaron, dibuja un recuadro con el nombre. Así el
% documento compila siempre.
% ---------------------------------------------------------------------
\newcommand{\figura}[2]{%
  \IfFileExists{figuras/#1}%
    {\includegraphics[width=#2]{#1}}%
    {\fbox{\parbox[c][0.45#2][c]{\dimexpr#2-2\fboxsep-2\fboxrule}{\centering\small\ttfamily figuras/#1\\(falta generar)}}}%
}

% =====================================================================
\title{\textbf{Simulación numérica del modelo de neuronas de Izhikevich:\\
tipos de disparo, error numérico y dinámica de red}}
\author{%
  [Nombre Apellido 1], [Nombre Apellido 2] y [Nombre Apellido 3]\\[2pt]
  \small Facultad de Matemática, Astronomía, Física y Computación, Universidad Nacional de Córdoba
}
\date{Redes Neuronales 2026 -- Trabajo Práctico 1}

\raggedbottom
\begin{document}

% El resumen ocupa todo el ancho, debajo del título. Las llaves {...} protegen
% los corchetes [ ] del texto, que si no cerrarían \twocolumn[...].
\twocolumn[{
  \begin{@twocolumnfalse}
    \maketitle
    \begin{abstract}
    % 5 a 10 líneas: objetivo + metodología numérica + resultados principales.
    Estudiamos numéricamente el modelo de neuronas de Izhikevich, un sistema
    bidimensional de ecuaciones diferenciales con reseteo que reproduce la variedad
    de patrones de disparo de neuronas corticales y talámicas con un costo
    computacional bajo. Integramos el modelo con el método de Runge--Kutta de
    orden~4 ($h=\num{0.1}$\,ms), aplicando el reseteo después de cada paso, y
    reproducimos los ocho paneles de la figura~2 de la Ref.~\cite{izhikevich2003}:
    regular spiking, intrinsically bursting, chattering, fast spiking,
    tálamo-cortical (dos protocolos), resonador y low-threshold spiking.
    Mostramos que el tratamiento discreto del reseteo reduce a orden~1 la
    convergencia de RK4 en los tiempos de disparo. [Si hacen la parte 2:]
    Finalmente, traducimos a Python el programa de la red de 1000 neuronas del
    artículo, reproducimos su figura~3 y mostramos que eliminar las
    autoconexiones de la matriz sináptica no modifica la dinámica colectiva.
    \end{abstract}
    \vspace{1em}
  \end{@twocolumnfalse}
}]

% =====================================================================
\section{Introducción}
% Qué es un modelo de neurona de disparo, por qué Izhikevich es relevante, citas, objetivo.
Las neuronas se comunican mediante potenciales de acción o \emph{spikes}:
excursiones bruscas y breves del potencial de membrana. Un modelo matemático
de neurona de disparo describe la evolución de ese potencial y reproduce la
generación de spikes ante una corriente de entrada. Estos modelos enfrentan un
compromiso entre realismo y eficiencia: el modelo de Hodgkin y
Huxley~\cite{hodgkin1952} describe el fenómeno con gran detalle biofísico, pero
con varias ecuaciones acopladas y decenas de parámetros, lo que lo vuelve
costoso para simular redes grandes; el modelo \emph{integrate-and-fire}, en
cambio, es muy eficiente pero no reproduce la riqueza de patrones de disparo
observada experimentalmente.

El modelo de Izhikevich~\cite{izhikevich2003,izhikevich2007} es un modelo
reducido de neurona de disparo que combina ambas propiedades: una dinámica
suficientemente rica como para reproducir diversos patrones de neuronas
biológicas y un costo computacional bajo. Su construcción está motivada por
reducciones, mediante teoría de bifurcaciones y formas normales, de modelos
biofísicamente más detallados de tipo Hodgkin--Huxley~\cite{izhikevich2007}.
Con sólo cuatro parámetros reproduce los principales tipos de disparo de
neuronas corticales y talámicas, y permite simular redes de miles de neuronas,
lo que lo convirtió en una herramienta estándar de la neurociencia
computacional.

El objetivo de este informe es integrar numéricamente el modelo, reproducir
sus tipos de disparo, analizar la precisión del método numérico
[y estudiar la dinámica colectiva de la red de 1000 neuronas del artículo]. % TODO

% =====================================================================
\section{Teoría}
\subsection{El modelo de Izhikevich}
Con la notación del artículo original~\cite{izhikevich2003}, $v$ en mV y $t$ en ms,
\begin{align}
  \dot v &= \num{0.04}\,v^2 + 5\,v + 140 - u + I(t), \label{eq:v}\\
  \dot u &= a\,(b\,v - u), \label{eq:u}
\end{align}
junto con la condición de reseteo
\begin{equation}
  \text{si } v \geq \SI{30}{\milivolt}: \qquad v \leftarrow c, \qquad u \leftarrow u + d .
  \label{eq:reset}
\end{equation}
La variable $v$ representa el potencial de membrana y $u$ es una variable de
recuperación que resume, de manera efectiva, procesos iónicos lentos (activación
de corrientes de K$^+$ e inactivación de corrientes de Na$^+$) y aporta
realimentación negativa sobre $v$. La entrada $I(t)$ representa una corriente
inyectada o una entrada sináptica efectiva. El parámetro $a$ controla la escala
temporal de recuperación de $u$ (valores chicos, recuperación lenta); $b$, cuánto
responde $u$ a las variaciones subumbrales de $v$; $c$ es el valor de reseteo de
$v$ después de un spike; y $d$, el incremento de $u$ después de cada spike.

El valor $v=\SI{30}{\milivolt}$ es el que usa el algoritmo para registrar el pico
del spike y aplicar el reseteo; \emph{no} debe interpretarse como un umbral
fisiológico fijo. En este modelo el umbral es dinámico: la iniciación del spike
depende del estado previo de $v$ y $u$. Por ejemplo, para $I$ constante el
reposo es la raíz menor de $\num{0.04}v^2+(5-b)v+140+I=0$, con $u^*=b\,v^*$
($v^*=\SI{-70}{\milivolt}$ para $b=\num{0.2}$ e $I=0$), y la raíz mayor es una
silla cuya variedad estable separa las trayectorias que disparan de las que
vuelven al reposo. [Si hay lugar: para $a=\num{0.02}$, $b=\num{0.2}$ el reposo
se pierde en $I_\mathrm{Hopf}\approx\num{3.8}$ e $I_\mathrm{SN}=4$.] % TODO

\subsection{Método numérico}
Con $\mathbf{x}=(v,u)$ y $\dot{\mathbf{x}}=\mathbf{f}(t,\mathbf{x})$, usamos
Runge--Kutta de orden~4:
\begin{equation}
  \mathbf{x}_{n+1} = \mathbf{x}_n + \tfrac{h}{6}\left(\mathbf{k}_1 + 2\mathbf{k}_2 + 2\mathbf{k}_3 + \mathbf{k}_4\right),
\end{equation}
con $\mathbf{k}_1=\mathbf{f}(t_n,\mathbf{x}_n)$,
$\mathbf{k}_{2,3}=\mathbf{f}(t_n+\tfrac h2,\mathbf{x}_n+\tfrac h2\mathbf{k}_{1,2})$ y
$\mathbf{k}_4=\mathbf{f}(t_n+h,\mathbf{x}_n+h\mathbf{k}_3)$, y $h=\num{0.1}$\,ms.
Después de cada paso se verifica si $v_{n+1}\geq\SI{30}{\milivolt}$; si se
cumple, se registra un spike en $t_{n+1}$ y se aplica inmediatamente el
reseteo~\eqref{eq:reset}. En todos los casos $u(0)=b\,v(0)$. En las figuras
el pico se dibuja en \SI{30}{\milivolt}.

% =====================================================================
\section{Resultados}
% Regla: la leyenda dice QUÉ se grafica; el texto dice QUÉ SE VE; el porqué va en Discusión.

\begin{figure}[t]
  \centering
  \figura{f1_base.png}{\columnwidth}
  \caption{Evolución temporal de $v(t)$ (arriba), $u(t)$ (medio) e $I(t)$
  (abajo) para $a=\num{0.02}$, $b=\num{0.2}$, $c=\SI{-65}{\milivolt}$, $d=2$ y
  un escalón $I=10$ desde $t=10$\,ms, integrados con RK4 ($h=\num{0.1}$\,ms)
  desde $v(0)=\SI{-70}{\milivolt}$, $u(0)=-14$. La línea discontinua marca el
  inicio del estímulo y los triángulos, los spikes registrados.}
  \label{fig:base}
\end{figure}

En la Figura~\ref{fig:base} se observa que, mientras $I=0$, la neurona permanece
en reposo en $\SI{-70}{\milivolt}$. Tras la aplicación del estímulo en $t=10$\,ms
emite el primer spike en $t=\num{13.5}$\,ms y otros seis hasta $t=100$\,ms. Los
intervalos entre spikes (ISI) crecen de \num{3.6} a \num{15.2}\,ms y desde
$t\approx45$\,ms el disparo es periódico con ISI $\approx19$\,ms
($\approx53$\,Hz). La variable $u$ salta $d=2$ en cada spike y relaja hacia
$b\,v$ entre spikes.

Para reproducir la figura~2 de la Ref.~\cite{izhikevich2003} usamos los
parámetros y condiciones iniciales de la Tabla~\ref{tab:parametros} y las
corrientes de la consigna: $I_1$, escalón de 10 en $t=20$\,ms; $I_2$, escalón de
1 en $t=20$\,ms; $I_3$, pulso de $-25$ en $[20,80)$\,ms; e $I_4$, \num{0.2} desde
$t=40$\,ms más un pulso de \num{2.2} en $[180,181)$\,ms.

\begin{table}[t]
  \centering
  \caption{Parámetros (notación de la Ref.~\cite{izhikevich2003}), corriente y
  condición inicial de cada panel. En todos los casos $u(0)=b\,v(0)$.}
  \label{tab:parametros}
  \small
  \begin{tabular}{lcccccc}
    \toprule
    Caso & $a$ & $b$ & $c$ [mV] & $d$ & $I$ & $v(0)$ [mV] \\
    \midrule
    RS  & \num{0.02} & \num{0.20} & $-65$ & 8          & $I_1$ & $-65$ \\
    IB  & \num{0.02} & \num{0.20} & $-55$ & 4          & $I_1$ & $-65$ \\
    CH  & \num{0.02} & \num{0.20} & $-50$ & 2          & $I_1$ & $-65$ \\
    FS  & \num{0.10} & \num{0.20} & $-65$ & 2          & $I_1$ & $-65$ \\
    TC1 & \num{0.02} & \num{0.25} & $-65$ & \num{0.05} & $I_2$ & $-63$ \\
    TC2 & \num{0.02} & \num{0.25} & $-65$ & \num{0.05} & $I_3$ & $-65$ \\
    RZ  & \num{0.10} & \num{0.26} & $-65$ & 2          & $I_4$ & $\num{-62.5}$ \\
    LTS & \num{0.02} & \num{0.25} & $-65$ & 2          & $I_1$ & $-65$ \\
    \bottomrule
  \end{tabular}
\end{table}

% figure* = figura a doble columna (va arriba de la página siguiente).
\begin{figure*}[t]
  \centering
  \figura{f2_figura2.png}{0.72\textwidth}
  \caption{Potencial de membrana $v(t)$ (azul) y corriente de entrada $I(t)$
  (negro, escala arbitraria) para las parametrizaciones RS, IB, CH, FS, TC1, TC2,
  RZ y LTS de la Tabla~\ref{tab:parametros}, integradas con RK4 y
  $h=\num{0.1}$\,ms. Duración: 200\,ms (300\,ms en RZ).}
  \label{fig:tipos}
\end{figure*}

En la Figura~\ref{fig:tipos} se observa que RS produce un tren de spikes cuyo
ISI crece de 21 a 45\,ms, mientras que FS dispara a $\approx133$\,Hz sin
adaptación apreciable. IB emite una ráfaga inicial de tres spikes seguida de
spikes simples cada $\approx31$\,ms, y CH dispara ráfagas de cinco a siete spikes
separadas por $\approx60$\,ms. LTS dispara a alta frecuencia con adaptación
notoria (ISI de \num{2.9} a \num{13.4}\,ms). En TC1 una corriente de sólo 1
alcanza para un disparo tónico lento, mientras que en TC2 la neurona se
hiperpolariza hasta $\approx\SI{-87}{\milivolt}$ durante el pulso negativo y, al
retirarlo, dispara una ráfaga de rebote de cinco spikes. En RZ la corriente
\num{0.2} produce oscilaciones subumbrales amortiguadas, sin spikes, y el pulso
de 1\,ms en $t=180$\,ms conmuta al sistema al disparo repetitivo (ISI
$\approx42$\,ms). Cualitativamente, las ocho curvas coinciden con las de la
figura~2 de la Ref.~\cite{izhikevich2003}. % TODO: verificar números

% Opcional, si hay lugar: figura de resonancia (figuras/f5_rz.png).
% El pulso de RZ no dispara si se aplica en t_p ∈ [73,85] o [123,130] ms
% (cerca de los mínimos de la oscilación subumbral) y sí en el resto.

\begin{figure}[t]
  \centering
  \figura{f6_error.png}{\columnwidth}
  \caption{Izquierda: máximo error absoluto en los tiempos de disparo de RS en
  función de $h$, para Euler, RK4 y RK4 con localización del cruce de
  \SI{30}{\milivolt} por interpolación lineal. Derecha: $v(t)$ con Euler
  ($h=\num{0.5}$\,ms) y RK4 ($h=\num{0.1}$\,ms); las líneas punteadas marcan los
  disparos de referencia (\texttt{solve\_ivp}, DOP853, tolerancia $10^{-11}$).}
  \label{fig:error}
\end{figure}

La Figura~\ref{fig:error} muestra que el error en los tiempos de disparo de
RK4 decrece linealmente con $h$, igual que el de Euler pero con una constante
menor: para $h=\num{0.1}$\,ms es de \num{0.52}\,ms contra \num{1.82}\,ms de
Euler. Al localizar el cruce de \SI{30}{\milivolt} dentro del paso el error baja
a \num{0.07}\,ms. % TODO

% ---- Segunda parte (opcional). Si no la hacen, borren esta subsección. ----
\subsection{Red de 1000 neuronas}
Tradujimos a Python el programa Matlab de la sección~IV de la
Ref.~\cite{izhikevich2003}: $N_e=800$ neuronas excitatorias con $a=\num{0.02}$,
$b=\num{0.2}$, $c=-65+15r_i^2$, $d=8-6r_i^2$ y $N_i=200$ inhibitorias con
$a=\num{0.02}+\num{0.08}r_i$, $b=\num{0.25}-\num{0.05}r_i$, $c=-65$, $d=2$
($r_i\sim U(0,1)$); matriz de pesos $S=[\,\num{0.5}R_E\mid -R_I\,]$; entrada
talámica $5\xi_i$ o $2\xi_i$ renovada cada 1\,ms; acoplamiento por eventos de
spike, $I_i^{syn}=\sum_{j\in F}s_{ij}$; y dos medios pasos de Euler de
\num{0.5}\,ms por milisegundo. Usamos la semilla 2026.

\begin{figure*}[t]
  \centering
  \figura{f8_autoconexiones.png}{0.85\textwidth}
  \caption{Red de 1000 neuronas durante 1000\,ms (semilla 2026). (a) Raster de
  la implementación literal del artículo, con la diagonal de $S$; (b) misma
  semilla con $s_{ii}=0$. Neuronas 1--800 excitatorias (negro) y 801--1000
  inhibitorias (rojo). Medio: tasa poblacional (bins de 1\,ms). Abajo:
  espectro de la tasa poblacional ($t>100$\,ms) y tasa media con y sin diagonal
  para 20 semillas.}
  \label{fig:red}
\end{figure*}

En la Figura~\ref{fig:red}(a) se observa actividad asincrónica de fondo
interrumpida por episodios de sincronización: bandas verticales en las que gran
parte de la población dispara junta, con un pico espectral cerca de 8\,Hz
(ritmo $\alpha$), como en la figura~3 de la Ref.~\cite{izhikevich2003}. La
red literal emite \num{7765} spikes (tasa media \num{7.77}\,Hz) y la red sin
autoconexiones \num{7561} (\num{7.56}\,Hz). Sobre 20 semillas, la diferencia de
tasa es $\num{0.02}\pm\num{0.12}$\,Hz. % TODO

% =====================================================================
\section{Discusión}
% ¿Se reproducen los tipos (incluida la resonancia)? Mecanismos. Fuentes de error. Paso h y reseteo.
Las simulaciones reproducen cualitativamente los ocho paneles de la
Ref.~\cite{izhikevich2003}: regular spiking, fast spiking, bursting (IB, CH),
low-threshold spiking, disparo tónico y de rebote tálamo-cortical, y la
resonancia de RZ. Las diferencias cuantitativas se deben a que el artículo no
tabula las corrientes, que aquí son las de la consigna. Los mecanismos se leen en
los parámetros: un $d$ grande produce adaptación (RS); un $c$ alto favorece las
ráfagas (IB, CH); un $a$ grande da recuperación rápida (FS); y un $b$ grande baja
el umbral (LTS, TC), produce el rebote tras la hiperpolarización (TC2) y vuelve
foco al reposo (RZ). En RZ el pulso dispara o no según la fase de la oscilación
subumbral en que llega: es una respuesta resonante, no de umbral fijo.

Como fuentes de error numérico destacamos: (i) el reseteo se detecta recién en
el primer punto de la grilla con $v\geq\SI{30}{\milivolt}$, lo que introduce un
error $\mathcal{O}(h)$ en cada tiempo de disparo, que se acumula de un spike al
siguiente y anula la ventaja del orden~4 de RK4 (Fig.~\ref{fig:error}); por eso
importan tanto el paso $h$ como el tratamiento de la condición de reseteo, y
localizar el cruce dentro del paso recupera precisión; (ii) la rigidez del
término $v^2$ cerca del spike, que con $h$ grande puede hacer divergir $v$ dentro
de un paso; y (iii) en casos cercanos a una bifurcación, como RZ, un error pequeño
puede cambiar el régimen de la trayectoria.

[Si hacen la parte 2:] Las autoconexiones aportan a cada neurona un único peso
(a lo sumo \num{0.5}\,mV o $-1$\,mV) frente a la suma de cientos de entradas:
quitarlas cambia los spikes individuales, porque la red es sensible a
perturbaciones, pero no la tasa media ni la estructura de los episodios
sincrónicos. [Opcional: Euler--Maruyama con ruido $\sigma\,dW$ y la variante con
$z(v)$, que es un modelo distinto y no reproduce la figura~3.] % TODO

% =====================================================================
\section*{Agradecimientos}
% Opcional.
Agradecemos a [los docentes de la materia] por [...].

% =====================================================================
% Referencias "a mano". Alternativa con BibTeX: ver la guía.
\begin{thebibliography}{9}\small
\bibitem{izhikevich2003}
  E.~M. Izhikevich, ``Simple model of spiking neurons,''
  \emph{IEEE Trans.\ Neural Netw.} \textbf{14}(6), 1569--1572 (2003).

\bibitem{izhikevich2007}
  E.~M. Izhikevich, \emph{Dynamical Systems in Neuroscience: The Geometry of
  Excitability and Bursting}, MIT Press, Cambridge, MA (2007).

\bibitem{hodgkin1952}
  A.~L. Hodgkin y A.~F. Huxley, ``A quantitative description of membrane
  current and its application to conduction and excitation in nerve,''
  \emph{J.\ Physiol.} \textbf{117}(4), 500--544 (1952).
\end{thebibliography}

\end{document}
